%\subsection{Costos fiscales del salario mínimo}
%\label{sec:costos-fiscales}

%El mínimo también determina obligaciones del Estado. El vínculo más directo es la pensión mínima, que aumenta con el piso salarial. Las pensiones superiores se reajustan, por regla general, con el IPC, respetando el piso de un mínimo. Cuando este crece más rápido, puede alcanzar mesadas que antes estaban por encima de él y ampliar el grupo sujeto a esa indexación. Por tanto, el costo persiste después del año del aumento y puede extenderse a más beneficiarios \citep{ley100}.

%El costo incluye también obligaciones con personas que todavía no se han pensionado. Un mínimo real más alto eleva el piso de las pensiones que se reconocerán en el futuro, siempre que el afiliado cumpla los requisitos de acceso. Además, puede adelantar el momento en que una mesada inicialmente superior al mínimo queda alcanzada por este. Por eso, evaluar solo el gasto del año siguiente deja fuera parte del efecto sobre los pagos futuros \citep{ley100,mfmp2026}.

%Como referencia de corto plazo, \citet{anifCosto2025} estimó un mayor gasto pensional del régimen de prima media de cerca de \$240.000 millones por cada punto porcentual adicional de aumento real del mínimo. Es una estimación publicada en noviembre de 2025 para la discusión del ajuste de 2026. No mide el valor presente de las obligaciones adicionales ni constituye una actualización para la negociación de 2027.

%El MFMP 2026 señala que cerca del 61\% de los pensionados de Colpensiones recibe un mínimo. También reconoce mayores cotizaciones, pues el 69,4\% de sus afiliados tiene una base de cotización de un mínimo. El balance fiscal requiere considerar ambos flujos. Su simulación de un punto porcentual adicional de crecimiento del mínimo cada año eleva las transferencias de la Nación a Colpensiones en 0,01 puntos del PIB en 2027 y 0,2 puntos en 2037 \citep{mfmp2026}.

%Ese ejercicio compara trayectorias acumulativas con una senda que ya incorpora un crecimiento del mínimo superior al IPC de 2,46\%. No mide el costo de pasar de 7\% a 8\% solamente en 2027, ni el ahorro de nuestra propuesta. Tampoco debe confundirse con indexar todas las pensiones al mínimo, un escenario distinto que presenta el MFMP \citep{mfmp2026}.


%El CARF calculó un mayor déficit de al menos \$5,3 billones en 2026 y \$8 billones desde 2027. Su estimación excluye respuestas del empleo y posibles compensaciones por demanda. Son cálculos de enero, anteriores al MFMP de junio, y no deben sumarse a sus proyecciones \citep{carf22}.

%Esas cifras anuales deben distinguirse del valor presente del costo fiscal. Este requiere proyectar y descontar el gasto adicional, neto de las cotizaciones adicionales, para los pensionados actuales y las futuras cohortes de pensionados. Un aumento real aplicado una sola vez modifica el nivel del mínimo de los años siguientes. Añadir un punto a su crecimiento cada año genera una trayectoria distinta. Las fuentes revisadas no proporcionan una sensibilidad del valor presente del pasivo pensional total a un punto de aumento real aislado.

%Las rentas vitalicias añaden una obligación distinta de las transferencias a Colpensiones. El mecanismo de cobertura del deslizamiento protege a las aseguradoras frente al riesgo de que el mínimo aumente más que la inflación bajo las condiciones de la regulación. Hacienda informó en diciembre de 2025 pagos acumulados cercanos a \$2,2 billones por ese mecanismo, de los cuales alrededor de \$1 billón correspondía a 2025. Esos pagos históricos no cuantifican el efecto del aumento de 2026 \citep{deslizamiento}.

%La nómina y los contratos requieren revisar las reglas de cada rubro. Que una remuneración no pueda ser inferior al mínimo no significa que toda la escala salarial pública deba subir en igual proporción. En los contratos de servicios, el impacto depende del componente laboral y de las condiciones de ajuste. Del mismo modo, un tope expresado en mínimos no implica por sí solo un desembolso presupuestal equivalente \citep{funcionPublica2026}.

%La recomendación de moderación incorpora así una restricción fiscal concreta. Añadir aumentos reales al piso salarial eleva compromisos pensionales que deben financiarse con cotizaciones, impuestos o deuda. El análisis debe conservar las compensaciones por mayores ingresos y distinguir obligaciones recurrentes, coberturas contingentes y gastos efectivamente pagados. Las fuentes consultadas no permiten calcular un ahorro fiscal total para nuestra propuesta de 2027.
